\documentclass[preview]{standalone}

\usepackage{amsmath}
\usepackage{amssymb}
\usepackage{stellar}
\usepackage{definitions}

\begin{document}

\id{psicologia-personalita-umanistica-sociale}
\genpage

\section{Teorie umanistiche e social-cognitive}

\begin{snippet}{teorie-umanistiche-introduzione}
    Gli approcci umanistici alla personalità si concentrano sulla totalità
    dell'esperienza personale conscia e sul potenziale di crescita
    dell'individuo. Valorizzano la personalità sana, la felicità e
    l'autorealizzazione.
\end{snippet}

\begin{snippetdefinition}{concetto-se-definition}{Concetto di sé}
    Il \emph{concetto di sé} è il modello mentale che un individuo possiede dei
    propri comportamenti tipici e delle proprie qualità personali.
\end{snippetdefinition}

\begin{snippetdefinition}{autorealizzazione-definition}{Autorealizzazione}
    L'\emph{autorealizzazione} è lo sforzo costante di realizzare il proprio
    potenziale e di sviluppare talenti e capacità. Maslow la colloca al vertice
    della gerarchia dei bisogni.
\end{snippetdefinition}

\begin{snippetdefinition}{accettazione-incondizionata-definition}{Accettazione incondizionata}
    L'\emph{accettazione incondizionata} è l'accettazione totale di un individuo
    da parte di un altro, senza imporre condizioni.
\end{snippetdefinition}

\begin{snippet}{rogers-horney}
    Per Rogers, la persona cerca di rendere il concetto di sé congruente con le
    esperienze reali di vita. Karen Horney sostenne invece che le persone
    possiedono un vero sé, che richiede condizioni ambientali favorevoli; in
    assenza di tali condizioni può svilupparsi un'ansia di base.
\end{snippet}

\begin{snippet}{caratteri-teorie-umanistiche}
    Le teorie umanistiche sono:
    \begin{itemize}
        \item \emph{olistiche}, perché interpretano gli atti individuali in
        rapporto alla personalità complessiva;
        \item \emph{disposizionali}, perché si concentrano sulle qualità innate
        che orientano il comportamento;
        \item \emph{fenomenologiche}, perché enfatizzano il punto di vista
        soggettivo dell'individuo.
    \end{itemize}
\end{snippet}

\begin{snippet}{valutazione-teorie-umanistiche}
    Le teorie umanistiche hanno valorizzato il potenziale di crescita della
    persona, ma sono state criticate per la vaghezza di alcuni concetti chiave,
    la difficoltà di verificarli empiricamente e la tendenza a trascurare alcune
    variabili ambientali.
\end{snippet}

\begin{snippet}{teorie-social-cognitive-introduzione}
    Le teorie cognitive e dell'apprendimento sociale cercano di collegare più
    direttamente personalità e comportamenti specifici. Considerano sia le
    circostanze ambientali sia i processi cognitivi e le differenze individuali.
\end{snippet}

\begin{snippetdefinition}{aspettativa-definition}{Aspettativa}
    L'\emph{aspettativa} è la credenza che un certo comportamento possa produrre
    una ricompensa in una determinata situazione.
\end{snippetdefinition}

\begin{snippetdefinition}{locus-of-control-definition}{Locus of control}
    Il \emph{locus of control}, o luogo di controllo, indica la credenza
    generale relativa al grado in cui i risultati ottenuti dipendono dalle
    proprie azioni oppure da fattori esterni.
\end{snippetdefinition}

\begin{snippet}{rotter-locus-control}
    Secondo Rotter, le aspettative derivano anche dalla storia personale dei
    rinforzi. Le persone con locus of control interno credono che i risultati
    dipendano soprattutto dalle proprie azioni; quelle con locus of control
    esterno credono che dipendano principalmente da fattori ambientali.
\end{snippet}

\begin{snippetdefinition}{determinismo-reciproco-definition}{Determinismo reciproco}
    Il \emph{determinismo reciproco} è il principio secondo cui individuo,
    comportamento e ambiente interagiscono reciprocamente: ciascuna componente
    influenza ed è influenzata dalle altre.
\end{snippetdefinition}

\begin{snippetdefinition}{autoefficacia-definition}{Autoefficacia}
    L'\emph{autoefficacia} è la credenza dell'individuo di poter eseguire una
    prestazione adeguata in una determinata situazione.
\end{snippetdefinition}

\begin{snippet}{bandura-autoefficacia}
    Per Bandura, il senso di autoefficacia influenza percezioni, motivazioni e
    prestazioni. Spesso non si tenta un compito quando ci si aspetta di essere
    inefficaci; anche con capacità e desiderio di agire, si può rinunciare se ci
    si ritiene non all'altezza.
\end{snippet}

\begin{snippet}{fonti-autoefficacia}
    Le principali fonti dei giudizi di autoefficacia sono:
    \begin{enumerate}
        \item risultati oggettivi;
        \item esperienza vicaria, cioè osservazione delle prestazioni altrui;
        \item persuasione da parte degli altri o di sé stessi;
        \item osservazione della propria attivazione emotiva durante il compito
        o mentre lo si immagina.
    \end{enumerate}
\end{snippet}

\includesnpt[width=55\%,src=/snippet/static-psychology/bandura-self-efficacy-model.jpg]{centered-img}

\end{document}
